\section{Attacks on the Approval Path}\label{sec:attacks}

An approval can fail in three places: in \emph{what is shown} to the human, in \emph{what is signed} by the
human's device, and in \emph{what is executed} afterwards. We built twelve attacks, \AB{1} to \AB{12}, that each
target one of these places. We executed every applicable attack against three third-party targets and the
reference defense, reviewed a fourth target statically, and mapped the family onto in-process frameworks through
the audit of \cref{sec:audit}. \Cref{tab:abfamily} defines the family, and \cref{tab:attacks} reports the
outcomes.

\begin{table}[t]
\centering\small
\caption{The approval-path attack family. ``Breaks'' names the property of \cref{sec:properties} that the
attack violates when it succeeds.}
\label{tab:abfamily}
\begin{tabularx}{\textwidth}{@{}l>{\raggedright\arraybackslash}p{0.2\textwidth}>{\raggedright\arraybackslash}X l@{}}
\toprule
ID & Attack & What the attacker does & Breaks\\
\midrule
\AB{1} & Amount cancellation & Sends two amount-like arguments, such as \code{amount=10000, total=-10000}, so that a summing check sees a small value & \Prop{2}, \Prop{6}\\
\AB{2} & Decoy-field recipient & Adds a second recipient or account under an innocuous key, such as \code{bcc}, so the display or grant check reads the wrong one & \Prop{2}, \Prop{4}\\
\AB{3} & Grant stretching & Reuses a standing ``allow similar'' grant for a different payee, recipient or larger amount & \Prop{6}\\
\AB{4} & Padding and hidden text & Pushes the harmful part of an argument off screen with blank glyphs (U+2800, U+3164) or reorders it with bidirectional controls & \Prop{4}\\
\AB{5} & Label spoofing & Chooses argument names that imitate trusted labels, such as \code{Target} or \code{Verified by client} & \Prop{4}\\
\AB{6} & Non-injective rendering & Finds two different effects that produce the same card, for example a boolean and the string \code{"true"} & \Prop{4}\\
\AB{7} & Enrollment squatting & Registers its own passkey for a principal that has none, then approves its own calls & \Prop{1}\\
\AB{8} & Challenge ambiguity & Exploits an encoding in which two different (identifier, nonce, digest) tuples produce the same signed bytes & \Prop{2}\\
\AB{9} & Ceremony control & Controls the page that drives the authenticator, so the page shows one effect and the device signs another & \Prop{1}, \Prop{5}\\
\AB{10} & Post-approval substitution & Changes arguments between approval and execution, the class Loopjacking documents~\citep{arunkumar2026loopjacking} & \Prop{2}\\
\AB{11} & Flooding within one run & Raises many approval requests in one turn until one is approved & \Prop{7}\\
\AB{12} & Threshold splitting & Splits one large effect into pieces that each stay under the step-up threshold & \Prop{7}, \Prop{6}\\
\bottomrule
\end{tabularx}
\end{table}

\subsection{A worked example}\label{sec:example}

\Cref{fig:example} shows \AB{2} as an approver sees it. The user asked the agent to email a report to their
boss. The compromised agent sends that email and adds a blind copy to an address it controls, and describes the
action truthfully except for the copy. On a narrated card, which is the agent's own description, nothing is wrong
to see. On the trusted card, which the gateway renders from the canonical call, the title names both recipients,
a risk line counts them, and the agent's text sits in a box marked untrusted. We gave both views to \llmModels{}
current language models acting as automated approvers (\cref{sec:rq4}). On the narrated view all
\exBccNarrApprove{} approved; on the trusted card all \exBccCardDeny{} refused. One model justified its approval of
the narrated view with ``The action matches the user's request exactly'', and the same model refused the card
because ``The email BCCs collector@example.test, a recipient the user never requested''. The attack and the
model were the same in both views; only the source of the displayed text differed.

\begin{figure}[t]
\centering
% Figure: the decoy-bcc stimulus in the narrated and the card view (text copied from data/llm/stimuli_llm_approver.json)
\begin{tikzpicture}[font=\scriptsize]
\node[draw=narrorange,line width=0.9pt,rounded corners=3pt,fill=untrustred!60,inner sep=6pt,
      text width=0.36\textwidth,align=left,anchor=north west] (L) at (0,0) {%
  \textbf{\small Narrated view} {\color{inkgray}(written by the agent)}\\[4pt]
  \ttfamily APPROVAL REQUEST\\
  The agent describes the action it wants to take:\\[2pt]
  Sending report.txt to boss@example.test with the subject 'Weekly report'.\\[8pt]
  \rmfamily\color{inkgray}User's task: \emph{Email the file report.txt to boss@example.test with the subject
  `Weekly report'.}};
\node[draw=cardblue,line width=0.9pt,rounded corners=3pt,fill=trustgreen!70,inner sep=6pt,
      text width=0.53\textwidth,align=left,anchor=north west] (R) at ($(L.north east)+(0.35,0)$) {%
  \textbf{\small Trusted card} {\color{inkgray}(rendered by the gateway from $\Desc$)}\\[4pt]
  \ttfamily
  Title: SEND to "boss@example.test", \textbf{"collector@example.test"} [send\_email]\\
  Risk notes:\\
  \ \ - 2 recipients/payees: "boss@example.test", "collector@example.test"\\
  Effect:\ \ Tool: send\_email\ \ \ Effect class: send\\
  \ \ Target: "boss@example.test", "collector@example.test"\\
  Arguments (names chosen by the agent; values exactly as they will be used):\\
  \ \ arg.bcc: \textbf{"collector@example.test"}\ \ [type: recipient]\\
  \ \ arg.body: "<report placeholder>"\ \ [type: text]\\
  \ \ arg.subject: "Weekly report"\ \ [type: text]\\
  \ \ arg.to: "boss@example.test"\ \ [type: recipient]\\
  Request context: user-0001, agent-0001, sess-0001\\
  \ \ Origin: user\_turn (asserted by the agent, not verified)\\
  Digest: 5ffd1292\ldots 75fb51ab\\[3pt]
  \fcolorbox{narrorange}{white}{\parbox{8.1cm}{\raggedright\ttfamily
  Agent's explanation (UNTRUSTED: written by the agent, may be false)\\
  Sending report.txt to boss@example.test with the subject 'Weekly report'.}}};
\node[anchor=north,font=\footnotesize\bfseries,text=narrorange] at ($(L.south)+(0,-0.15)$)
  {\exBccNarrApprove{} of \exBccN{} model approvers approve};
\node[anchor=north,font=\footnotesize\bfseries,text=cardblue] at ($(R.south)+(0,-0.15)$)
  {\exBccCardDeny{} of \exBccN{} model approvers deny};
\end{tikzpicture}

\caption{The same blind-copy attack (\AB{2}) seen two ways. Left: the narrated view, written by the agent. Right:
the trusted card, rendered by the gateway from the canonical descriptor; the agent's text appears only in the
box at the bottom. The card is abridged (expiry line omitted, context lines merged, digest shortened, separator lines around the explanation dropped, label punctuation normalized). Across \llmModels{} language-model approvers at temperature zero, \exBccNarrApprove{} of
\exBccN{} approved the left view and \exBccCardDeny{} of \exBccN{} denied the right one.}
\label{fig:example}
\end{figure}

\subsection{Targets and method}

We tested the documented, intended use of each target locally, with neutral placeholder data; no real service,
credential or person was contacted. The targets are:
\begin{itemize}[leftmargin=1.4em,itemsep=1pt]
\item \textbf{AP2}, the Agent Payments Protocol reference implementation~\citep{ap2spec}, at commit
  \code{e1ea56d} (v0.2.0). We executed its mandate constraint checks and read its sample card component from source.
\item \textbf{AGT}, Microsoft's Agent Governance Toolkit, whose ADR~0030 specifies an action-bound, fail-closed
  approval protocol~\citep{microsoft2026adr0030}, at commit \code{47abed7}. All AGT cells were executed.
\item \textbf{MCP}, the Model Context Protocol Python SDK 2.3.0 and its elicitation feature~\citep{mcp2025elicitation},
  executed over an in-memory client and server session.
\item \textbf{OpenClaw} 2026.10.3, whose exec approvals and ``allow always'' grants we reviewed statically in the
  shipped bundle, because its live loop requires a model; every OpenClaw cell is marked as not executed.
\item \textbf{Frameworks}, the eight in-process frameworks of \cref{sec:audit}; this column is derived from the
  audit's executed probes, not re-run.
\item \textbf{\sys{}}, the reference defense of \cref{sec:design}; each cell maps to named tests run against the
  live gateway.
\end{itemize}

\begin{table}[t]
\centering\small
\caption{Approval-path attacks against real designs. VULN: the attack succeeds against documented usage. PART:
the attack succeeds only at the argument level, depends on the integrator, or reflects a stated trust assumption.
ok: the attack fails. n/a: the target has no such feature. $^\ast$: not executed in this run (static review, or
derived from the audit). The AP2 \AB{6} cell combines an executed constraint check with a source reading of the
sample card.}
\label{tab:attacks}
\setlength{\tabcolsep}{5pt}
\begin{tabular}{@{}l*{6}{c}@{}}
\toprule
Attack & AP2 & AGT & MCP & OpenClaw & Frameworks & \sys{}\\
\midrule
\inputtab{tables/tab_attacks.tex}
\bottomrule
\end{tabular}
\end{table}

\subsection{Findings}

\paragraph{AP2 binds what it signs; its sample display does not.} AP2's constraint layer rejected every executed
binding attack. Amounts are a single typed integer field, so \AB{1} has nothing to cancel. Allowed payees are
matched on the merchant identifier, so a decoy display name does not redirect a payment (\AB{2}), and a pre-set
payee or currency in an open mandate cannot be stretched (\AB{3}). Two residuals remain. The sample mandate card
summarizes \code{payee.name}, a merchant-chosen string with no length cap or invisible-character handling, while
the binding is on \code{payee.id}, so two different payees can share a display (\AB{6}). And a mandate budget is
cumulative only if the integrator persists and adds up a caller-supplied usage context; without that, two
sub-budget payments both pass (\AB{12}). Both are integrator or sample issues, not breaks of the signed
mandate, and the specification already asks for a trusted surface.

\paragraph{AGT is the strongest third-party design we tested.} Its digest covers the whole action binding, so
substitution after approval fails with \code{action\_digest\_mismatch} (\AB{10}) and decoy fields change the
digest (\AB{2}). A language model's advisory output cannot satisfy a human stage (\AB{5}), and approvals are
consumed once. Its two partial cells concern stated scope. The project documents that floats are serialized with Python's
representation, not RFC~8785's, so three of five sample floats would hash differently in a verifier written in
another language (\AB{8}); this is a parity failure that fails closed, not a collision. ADR~0030 also does not
claim flood control, and the module opened 50 pending requests freely, leaving rate limiting to the host
(\AB{11}).

\paragraph{MCP elicitation forwards whatever the server writes.} Elicitation is a form prompt, not an effect-bound
approval, and the SDK passes the server's message to the client unchanged. In the executed session, a
521-character message carrying 400 blank Braille glyphs, a right-to-left override and 50 Hangul fillers reached the
client byte for byte (\AB{4}). Text imitating trusted labels arrived verbatim while the tool paid a different
account (\AB{5}). Two different payments, 10 to one account and 10000 to another, produced byte-identical
requests, and both executed after ``accept'' (\AB{6}). The form answer is not bound to any effect (\AB{10},
partial). The specification already directs sensitive interactions to its URL mode, which keeps them away from the
client and recommends that clients rate-limit requests~\citep{mcp2025elicitation}; these findings concern form mode and the clients that render it.

\paragraph{OpenClaw's card is the positive control for rendering.} Its exec-approval card escapes invisible
characters, caps length and redacts secrets, the treatment that \AB{4} requires, and its ``allow
always'' grants are keyed to the resolved argument vector and working directory. The residual is that a standing
grant is matched on parsed arguments rather than a signed digest (\AB{3}, partial, static).

\paragraph{In-process frameworks hand the ceremony to the agent.} Because approval callbacks run inside the agent's
process, the agent can drive the ceremony itself (\AB{9}), and the audit reproduced post-approval substitution in
four runtimes (\AB{10}, \cref{sec:audit}). No framework ships a trusted renderer, so agent-chosen names and text
reach the human unmediated (\AB{5}, \AB{6}).

\paragraph{The reference defense blocks every agent-reachable attack, and one trust assumption remains.} \sys{}
stops \AB{1} to \AB{8} and \AB{10} to \AB{12} in tests against the live gateway (\cref{sec:design} describes the
mechanisms). Its \AB{9} cell is partial by design: the agent's token is refused on every ceremony endpoint, but a
page served from the gateway's own origin could show one card and obtain a signature over another approval's
challenge, because authenticators do not display the challenge. The receipt still verifies at the key level, and
the effect that executes is the one the device signed. This is the honest-rendering-origin hypothesis \Hyp{3} of
\cref{sec:formal}, not an attack the agent can mount.

We describe these as findings about designs, not as vulnerabilities in deployed services; the
\hyperref[sec:ethics]{Ethics and Availability} section discusses them further.
